\chapter
[\CO{[EKG]} EKG]
{\CC{[EKG]}\\EKG}
\label{chp:EKG}

\minitoc

\section{Grundlagen}

Erregung Endokard → Epikard

Repolaristaion umgekehrt: Epikard → endokard → pos. p-Welle

1 Kasterl: 40 ms

5 Kasterl: 200 ms

Amplitude  enspricht in etwa der Muskelmasse

\subsection{Reizleitungssystem udn Reizbildung}


\ce{NA+} Depolarisation

\ce{K+} Repolarisation

Wenn Brustwandableitungen zu hoch geklebt werden, erhält man Ischämie-Zeichen im EKG.

\VonBis{$V_{1R}$}{$V_{6R}$} Spiegelbild

$V_{1R}$ = $V_2$

$V_{2R}$ = $V_1$

Ableitungen nach Nehb: Dorsal (D), Inferior (I), Anterior (A)

p \VonBis{50}{100}

QRS \VonBis{60}{100}

\VonBis{101}{119} inkompletter Schenkelblock

\VonBis{120}{\ldots} kompletter Schenkelblock

Q, S \ldots{} nicht obligat

QTc \ldots{} Q (R) bis Ende T

<390 Männer, <440 Frauen

\section{Beurteilung des EKGs}

\begin{itemize}
  \item Rhythmus
  \item Frequenz
  \item Lage
  \item Rotation
  \item P-Welle
  \item PQ
  \item QRS-Komplex
  \item ST
  \item T-, U-Welle
  \item QT
\end{itemize}

\subsection{"`Anatomie der EKG-Kurve"'}

\subsubsection{p}

Normalbefund:

in I (\& II) positiv

kann in III, $V_1$ negativ sein

PQ konstant

QRS auf p

normfrequent

max 2\nicefrac{1}{2} Kästchen


Typische pathologische Befunde:

\begin{itemize}
  \item Junktionaler Ersatzrhythmus
  \item Artefakte
  \item Absolute Arrhythmie bei Vorhofflimmern
  \item P-Wellen-Asystolie
  \item Torsades des pointes
\end{itemize}

hohes p: dectroarteriale: Cor pulmonale, pulm. Hypertonie, hereditäres Vitium

breites p: p sinoatriale/mitrale, li VH, $V_1$ biphasisch

breites und hohes p: p bicordiale

neg. p in II, III: retrograde Erregung; in I: Situs inversus

Ursachen: 
geschädigtes VH-Myokard, ektope Erregunf, retrograde VH-Erregung, evtl. in QRS versteckt, Flimmer-/Flatterwellen





multifokale atriale Tachykardie: Lungenerkrankung, Elyte, Digitalis, metabolisch. Rx: causal, Ca-Antagonisten

???-junktion. Ersatzrhythmus (Katecholamine?)


\subsubsection{PQ}

< 120

< WPW, LGL, Basaler VH-Rhythmus (oberer Knotenrhythmus)

> 200 

Normal, AV-Block. Vagotonie, KHK, CMP, Digitalis

Erkenne: AV-Block \VonBis{I}{III}\textdegree

\subsubsection{Lagetyp}

\Ca{Der überdrehte Rechtstyp un der Rechtstyp sind beim Erwachsenen immer pathologisch!} Steiltyp altersabhängig (normal bei jung und schlank), überdrehter Linkstyp i,\,d.\,R. pathologisch

erkenne: verpoltes EKG

\subsubsection{Rotation}

rechts: Thorax: Emphysem-Th., Kyphoskoliose, Re-Herz-Hypertrophie, Seitenwandinfarkt, Infarktnarbe, linksposteriorer Hemiblock

links: Adipositas, li-Herz-Belastung, Hinterwandinfarkt/Narbe, Linksanteriorer Hemiblock

\subsubsection{QRS}

\subsubsection{Q}

\subsubsection{S}

\subsubsection{ST}

\subsubsection{T}

\subsubsection{U}

\subsubsection{QT, QTc}

\section{Veränderungen in der EKG-Kurve}

\subsection{Hypertrophie}

\subsection{Störung der Reizleitung}

\subsubsection{AV-Block}

\subsubsection{Schenkelblock}

\Cave{Ischämiezeichen sind bei Schenkelblock und Ischämiezeichen \EE{nicht anwendbar}!}

\subsection{Ischmämiezeichen}



